\documentclass[12pt]{beamer}

\usepackage[T2A]{fontenc}        
\usepackage[utf8]{inputenc}  
\usepackage[russian]{babel}

\usepackage{multirow} 
\usepackage{graphicx}

\usepackage{amsmath}
\usepackage{amssymb} 

\usepackage{hyperref}
\title{Презентация}
\date{\today}

\begin{document}
\begin{frame}
\titlepage
\end{frame}

\section{Формулы}
\begin{frame}{Пределы}
\begin{equation}
\int \sqrt{a^2-x^2} \, dx = \frac{x}{2}\sqrt{a^2-x^2}+ \frac{a^2}{2}\arcsin \frac{x}{a} +C
\end{equation}
\begin{equation}
\lim_{n \to \infty} \frac{ x\tan^{2n}\frac{\pi x}{4}}{\tan^{2n}\frac{\pi x}{4} + 1}
\end{equation}
\end{frame}

\begin{frame}{Интегралы}
\begin{equation}
\iint\limits_{\Omega} f(\sqrt{x^2+y^2}) \,dx\,dy , \quad где \Omega =\{|y| \le |x| ; |x| \le 1\}
\end{equation}

\begin{equation}
\int\limits_{e}^{+\infty} \dfrac{1}{x^p \left(\ln x\right)^q \left(\ln\ln x\right)^r} \,dx
\end{equation}
\end{frame}

\begin{frame}{Ряды}
\begin{equation}
f(x)=\sum\limits_{n=1}^{\infty} \dfrac{x}{(1+x^2)^n}
\end{equation}

\begin{equation}
\sum\limits_{n=1}^{\infty} \left(-1\right)^{n-1} \frac{2^n \sin^{2n} x}{2}
\end{equation}
\end{frame}

\section{Таблица}
\begin{frame}{Таблица}
\begin{table}[h]
\centering
\begin{tabular}{|c|c|c|}
\hline
\multicolumn{3}{|c|}{Заголовок} \\
\hline
\multirow{2}{*}{A} & B & C \\
\cline{2-3}
& D & E \\
\hline
F & G & H \\
\hline
\end{tabular}
\end{table}
\end{frame}

\section{Списки}

\begin{frame}{Маркированный список}
\begin{itemize}
\item Пределы
\item Интегралы
\item Ряды
\end{itemize}
\end{frame}

\begin{frame}{Нумерованный список}
\begin{enumerate}
\item Пределы
\item Интегралы
\item Ряды
\end{enumerate}
\end{frame}

\section{Литература}
\begin{frame}{Список литературы}
\begin{thebibliography}{99}
\bibitem{art1} \href{https://kvm.gubkin.ru/pub/uok/Demidovich.pdf}{Демидович Б.\,П. Сборник задач и упражнений по математическому анализу.}
\bibitem{art2} А.\,П. Комбаров, Ю.\,В. Садовничий. Аналитическая геометрия. 
\end{thebibliography}
\end{frame}
\end{document}